\section{The same arithmetic discrepancy in the mixed filtration and the source boundary}
\label{CFStage:sec:mixed-amplification-identity}

The following calculation connects the independent mixed faces, the
nilpotent filtration and the original phase observation on the same
arithmetic class. All maps used here are constructed in the preceding
attachments; this section computes their composite on the full packet.

Fix an actual zero $\rho=1/2+\delta+i\gamma$ in the open strip,
with $0<\delta<1/2$ and $\gamma>0$, its full order $m$, and a
full packet $h$ containing it. The other offcritical quadrants
map to this quadrant by the original reflection and conjugation
arrows of (CAU.30)--(CAU.32), retaining their full local jets.
Retain
the original CRT idempotent and terminal jet
\[
 v_\rho=e_\rho(s-\rho)^{m-1}\in E_h,
 \qquad u_\rho=\sigma_hv_\rho\in T_UQ,
 \qquad 0<\delta<1/2,\quad\gamma>0.
 \tag{MAI1}
\]
The coefficient of $(s-\rho)^{m-1}$ in this local jet is one;
therefore $v_\rho\ne0$. The full-unit section satisfies
$J_hs_h=1$, so $\sigma_h$ is injective and $u_\rho\ne0$.
The nilpotent kills $v_\rho$, and the full original exponential
therefore gives
\[
 A_av_\rho=a^\rho v_\rho,
 \qquad \overline{\mathcal R}_a u_\rho=a^\rho u_\rho
 \quad(a>0).
 \tag{MAI2}
\]
The second identity uses the actual theta-boundary homotopy, not an
equivariance assertion for the section before quotienting.

For $n\geq1$, in the coefficient algebra $B_{h,n}=E_h[t]/(t^n+1)$, the original
insertion and trace retraction satisfy
\[
 E_h\xrightarrow{\iota_n}B_{h,n}\xrightarrow{r_n}E_h,
 \qquad r_n\iota_n=1,
 \qquad \Psi_{q_0}\iota_n=\iota_n,
 \qquad r_n\Psi_{q_0}=r_n.
 \tag{MAI3}
\]
Here $q_0$ is an odd prime power coprime to $n$, and $r_n$ is $1/n$ times
the trace of multiplication. The complete coefficient computation
is $r_n(\sum b_it^i)=b_0$. Thus $\iota_nv_\rho$ is nonzero,
fixed by coefficient Frobenius, and has arithmetic eigenvalue
$a^\rho$. It is represented on the unsynchronized face $\{0\}$
by its supported constant coefficient and absent remaining slots.
Synchronization also observes it on the full face. Both vertices
and their connecting arrow are retained; (MAI3) identifies their
unchanged arithmetic amplitude.

Tensor each of these maps over $\mathbb C$ on its amplitude
spaces. To prove nonvanishing explicitly, choose a linear
functional $\ell$ with $\ell(v_\rho)=1$. Then
$\ell^{\otimes k}(v_\rho^{\otimes k})=1$; similarly the tensor
retraction $r_n^{\otimes k}$ carries
$(\iota_nv_\rho)^{\otimes k}$ to $v_\rho^{\otimes k}$.
The injective original cohomology map carries the latter to
$u_\rho^{\otimes k}\ne0$. These tensor arrows retain the ordered
tuple of coefficient-face labels. Their arithmetic eigenvalue is
$a^{k\rho}$, and hence
\[
 \log\!\left(\frac{|a^{k\rho}|^2}{a^k}\right)
       =2k\delta\log a.
 \tag{MAI4}
\]

Now retain the entire tensor of the original length-$m$ local
blocks, not only this terminal line. In the original jet basis
write $N_{\rm tot}=\sum_{\nu=1}^kN_\nu$ and use the filtration
proved in (MW15). On every nonzero graded piece $\operatorname{Gr}_j^M$ the added
Tate-scaling operator $F_q^{\otimes k}$, for real $q>1$, acts by
$q^{k\rho+j/2}$. Indeed on a basis tensor of indices
$j_1,\ldots,j_k$ its scalar is the product
$\prod_\nu q^{\rho+j_\nu/2}$, and their sum is $j$.
The coefficient operator on a retained eigenspace contributes a
root of unity $\omega$. Every eigenvalue on that graded
eigenspace is therefore
\[
 \lambda_{k,j,\omega}=\omega q^{k\rho+j/2},
 \qquad
 2\log_q|\lambda_{k,j,\omega}|-(k+j)=2k\delta.
 \tag{MAI5}
\]
The equality uses $|\omega|=1$ and the original real logarithm
of $q$. It holds for every chosen nonzero graded piece, including $j=0$.
For the tensor of the terminal lines in (MAI1), the index is
exactly $j=-k(m-1)$ and $\omega=1$. Thus this computation
includes a specified nonzero line in the arithmetic retract.

The exact reconstruction from (MW19) also tensors without
commuting operators in an unproved order:
\[
 A_q^{\otimes k}
 =F_q^{\otimes k}(\Gamma_q^{-1})^{\otimes k}
                 \exp((\log q)N_{\rm tot}).
 \tag{MAI6}
\]
Each tensor factor has that ordered identity. Operators in
different factors commute, so multiplying their finite
exponentials gives the displayed sum in the exponent. On the
graded piece, the exponential induces the identity and
$(\Gamma_q^{-1})^{\otimes k}$ contributes $q^{-j/2}$.
The resulting arithmetic eigenvalue is exactly $q^{k\rho}$.
This proves the map between the added Tate-scaling observation
and the original arithmetic action, including every nilpotent
term. Integer monodromy indices alone do not change its real
exponent.

For its original source-norm observation, take the reflected
quartet packet $h_0$ of (PAM.1), with $\delta>0$, and the complete
cyclic tensor factor of (PAM.2)--(PAM.5). Its nonzero terminal
eigenvector $v_k$ has eigenvalue $k\rho$. Its exact relation to
the tensor in (MAI4) includes a nonzero scalar. Put $K=k(m-1)$.
The extreme tensor centre $k\rho$ can arise only by choosing
$\rho$ in every factor of that quartet. On this block, the
multinomial expansion gives
$(N_1+\cdots+N_k)^K=K!/((m-1)!)^k\prod_\nu N_\nu^{m-1}$.
Every other term vanishes because some exponent is at least $m$.
The complete unit of (PAM.3) acts on this terminal monomial by
its nonzero scalar value. Consequently
\[
 \eta_kv_k=c_{h_0,k}v_\rho^{\otimes k},\qquad
 \mathcal I_kv_k=c_{h_0,k}u_\rho^{\otimes k},\qquad
 c_{h_0,k}=\frac{(k(m-1))!}{((m-1)!)^k}
                    \bigl((g/h_0)(\rho)\bigr)^k\ne0.
 \tag{MAI6a}
\]
Here $v_\rho$ in the first equality is taken in $E_{h_0}$.
If the initially chosen $h$ differs from $h_0$, use the common
full packet $H=\operatorname{lcm}(h,h_0)$. Both terminal jets
have the same image in $E_H$, and the literal identities
$\sigma_H\iota_{hH}=\sigma_h$ and
$\sigma_H\iota_{h_0H}=\sigma_{h_0}$ identify their original
arithmetic classes. This proves the second equality with the
original $u_\rho$ of (MAI1), retaining the entire coefficient.
For this phase observation take $k\geq2$, $L>0$,
$0\leq\theta<2\pi$, and an admitted polynomial cutoff
$N\geq q_k-1$, with $q_k=(1+k(m-1))(k+1)^2$.
Write $R_k=R_{N,\theta}$ for the canonical lift on this specified
source of (PAM.9)--(PAM.13). The source lift
$R_kv_k$ is nonzero because its original cohomology observation
is injective. With the exact twisted phase domain and its
generator $D_{L,\theta}=-\partial_r+k/2$, let
\[
 B_k=D_{L,\theta}R_k-R_kA_k.
 \tag{MAI7}
\]
Every term belongs to its proved domain. Put $f=R_kv_k$.
The spectral expansion of the skew-adjoint derivative on that
domain gives frequencies $u_n=(2\pi n+\theta)/L$ and
\[
 \frac{\|B_kv_k\|^2}{\|R_kv_k\|^2}
 =k^2\delta^2+
   \frac{\sum_n(u_n-k\gamma)^2\|f_n\|^2}
        {\sum_n\|f_n\|^2}
 \ \geq\ k^2\delta^2.
 \tag{MAI8}
\]
To verify the equality, apply $D_{L,\theta}-k\rho$ to each
orthogonal frequency component. Its multiplier is
$-k\delta+i(u_n-k\gamma)$, whose squared modulus is
$k^2\delta^2+(u_n-k\gamma)^2$. Parseval and the domain
condition justify both sums. The denominator is positive
because $f\ne0$.

Equations (MAI4), (MAI5) and (MAI8) are thus connected
observations of the same original arithmetic displacement.
The first follows through the coefficient-pure retract; the
second computes it after the complete monodromy grading; the
third measures its necessary original theta-boundary energy.
In particular the fully evaluated phase calculation does not
erase it: at the first cutoff the boundary is the nonzero
residue map of (PSC.50)--(PSC.64), even when the two section
gluing maps agree. The arithmetic four-window identities in
(AAM.19)--(AAM.30) retain that same boundary scale and all its
inverse Gram factors.

Deligne's final strict-weight exclusion in Weil II applies to
geometric constituents whose integer weights and strict upper
bound have both been proved: the complete calculation is
recorded in the following dyadic-bootstrap chapter. In (MAI5)
the integer $j$ has been computed, whereas the central real
weight is the unchanged $2\Re\rho$. Calling that central weight
one would replace the explicitly computed scalar rather than
deduce its value. The present calculation identifies its exact
location in the assembled mixed object and in the original
boundary operator; it proves no disappearance of the marked
cohomology $T_UQ$.
